\section{The actual compact-unitary automorphic decomposition}
\label{sec:automorphic-decomposition-v43}

The preceding classical-parts theorem preserves exact level and ordinary
conditions inside a supplied decomposition.  Here we construct the
underlying automorphic decomposition on the original definite unitary
quotient.  The endpoints are the full original \texttt{formU} and
\texttt{decompU}.  These endpoints are independent of the arithmetic
construction of a Galois representation for an automorphic constituent.
The latter construction remains necessary to obtain a full
\texttt{PartAltAlt} and \texttt{AutomorphicHeckeSet}.

Fix the original \(\ell,L,s,n,V,Q,c,\tau_0,\iota\), with
\(n\ge1\).  Write \(G=U_{\mathbb A}\),
\(\Gamma=U(L^+)\), \(G_f=U_f\), and \(G_\infty=U_\infty\),
with the literal identity Hermitian form and embeddings of FNL2.  The
source \(\mathcal A=\mathrm{autU}(L,n)\) consists of continuous,
left \(\Gamma\)-invariant complex functions on \(G\), smooth under
\(G_f\), and finite under \(G_\infty\).  Its action is right
translation.  An original automorphic representation is a nonzero
irreducible invariant algebraic subspace of this very function space.

For the specified weight and basis, let
\[
 D=\bigotimes_{a\in\operatorname{Emb}_\ell}
             \nabla_{\lambda(a)}(\mathbf C),\qquad
 E:\mathbf C^d\xrightarrow{\sim}D,
\]
and write \(R=\mathrm{infRep}(\iota,\lambda)\).  The letter \(E\)
is the exact source basis equivalence, not a newly chosen realization.
Let
\[
 p:G\longrightarrow G^a\times G^\ell,
 \qquad
 p(g)=\bigl(\mathrm{away}(g),(\mathrm{cLoc}_{w(v)}(g))_{v\mid\ell}\bigr).
\]
For an algebraic coefficient function \(F\), the exact source scalarization is
\begin{equation}\label{eq:adp-toform}
 \mathcal T_F(\lambda^*)(g)
       =\lambda^*\bigl(R(g^{-1})E(F(p(g)))\bigr),\qquad
                  \lambda^*\in D^*.
\end{equation}
The symbols \(\lambda\) and \(\lambda^*\) respectively denote the fixed
weight and an arbitrary linear functional; there is no dual-weight change
implicit in this notation.

\subsection{Algebraic semisimplicity on the original quotient}

% BEGIN ADP1
\paragraph{ADP1. A compact fundamental set and the actual positive invariant form.}
\label{claim:ADP1}
The quotient \(\Gamma\backslash G\) is compact and has a positive,
finite, right-invariant measure.  On \(\mathcal A\) the formula
\[
       \langle f,h\rangle
       =\int_{\Gamma\backslash G}\overline{f(g)}h(g)\,d\mu(g)
\]
is a positive definite, right-invariant Hermitian form.

Indeed choose an open compact subgroup \(K_f\le G_f\).  Such subgroups
are obtained by imposing integral matrices at almost all places and open
compact congruence conditions at the finitely many remaining places; the
unitary equations cut out closed subgroups.  FNL3 gives finitely many
\(t_i\) such that every \(g\) is \(\gamma t_i z k\), with
\(z\in G_\infty\), \(k\in K_f\).  The group \(G_\infty\) is a
finite product of compact identity-form unitary groups: at every real
place of \(L^+\), the two complex embeddings of \(L\) are conjugate,
and the equation is \(g^*g=1\).  Thus
\(\bigcup_it_iG_\infty K_f\) is a compact fundamental set.

The rational subgroup is discrete in these same adelic coordinates,
and FV5--FV8 give unimodularity and the invariant quotient measure for
this unitary group.  Compactness makes its measure finite; a nonempty
open subset has positive measure by the local Haar quotient formula.
Continuous functions on the compact quotient are bounded, so the
integral is defined.  Right invariance proves invariance of the form.
If the squared norm is zero and \(f(g_0)\ne0\), continuity supplies
a neighborhood of \(g_0\) on which \(|f|\) is bounded below by a
positive constant.  Its image in the quotient has positive measure,
contradicting the zero integral.  Hence the form is definite on the
actual functions, without passing to almost-everywhere equivalence classes.
% END ADP1

% BEGIN ADP2
\paragraph{ADP2. Finite multiplicity from the exact finite class representatives.}
\label{claim:ADP2}
Let \(S\) be a finite-dimensional complex representation of
\(G_\infty\), and let
\[
 \mathcal E_S=\operatorname{Hom}_{G_\infty}(S,\mathcal A).
\]
For every open \(K_f\le G_f\), the space \(\mathcal E_S^{K_f}\)
is finite-dimensional.  If \(t_1,\ldots,t_r\) are the representatives
from FNL3, the evaluation map
\[
 \mathcal E_S^{K_f}\longrightarrow (S^*)^r,
       \quad T\longmapsto\bigl(v\mapsto T(v)(t_i)\bigr)_i
\]
is injective.

If all evaluations vanish, write \(g=\gamma t_i z k\) in the exact
FNL3 order.  Left rational invariance removes \(\gamma\); right
\(K_f\)-invariance removes \(k\); equivariance gives
\(T(v)(t_i z)=T(zv)(t_i)=0\).  Thus \(T=0\), proving the bound
\(\dim\mathcal E_S^{K_f}\le r\dim S\).  Neither independence of the
representatives nor triviality of their stabilizers is required.

Each \(T\in\mathcal E_S\) is fixed by an open compact subgroup of
\(G_f\): take a finite basis of \(S\), intersect the open finite
stabilizers of its images, and choose an open compact subgroup inside
that intersection.  Every compact subgroup has finite orbits on each
element of \(\mathcal E_S\), since an open stabilizer has finite
index in a compact group.  These are the exact admissibility and
finite-orbit assertions needed below.
% END ADP2

% BEGIN ADP3
\paragraph{ADP3. Complete reducibility in each actual multiplicity space.}
\label{claim:ADP3}
Each \(\mathcal E_S\) is an algebraic sum of irreducible
\(G_f\)-subrepresentations.

Choose any basis \(s_1,\ldots,s_d\) of \(S\) and give
\(\mathcal E_S\) the Hermitian form
\[
       (T,T')_S=\sum_{j=1}^d
                         \langle T(s_j),T'(s_j)\rangle.
\]
ADP1 makes it positive definite and \(G_f\)-invariant.  For an arbitrary
\(T\), ADP2 gives an open compact \(K_f\) fixing it, with finite
dimensional fixed space and finite orbits on all other vectors.  The
abstract argument in PKT5--PKT6 now applies with precisely this
\(G_f,\mathcal E_S,K_f\) and this Hermitian form.  It expresses
\(T\) as a finite sum of vectors in irreducible invariant subspaces.
All hypotheses of that argument have just been produced here; none of
the separate geometric multiplicity inputs used in PKT4 are imported.
% END ADP3

% BEGIN ADP4
\paragraph{ADP4. Algebraic decomposition of the complete automorphic function space.}
\label{claim:ADP4}
Every \(f\in\mathcal A\) belongs to the algebraic sum of the original
automorphic representations in \(\mathcal A\).

Its \(G_\infty\)-orbit spans a finite-dimensional invariant space by
the definition of \(\mathcal A\).  The form of ADP1 makes this space
unitary, so successive invariant orthogonal complements decompose it
into finite-dimensional irreducibles \(S\).  For each such \(S\),
ADP3 decomposes its multiplicity space \(\mathcal E_S\) under \(G_f\).
It remains to check that, for an irreducible subspace
\(U\subset\mathcal E_S\), its evaluation image
\[
       \operatorname{ev}(S\otimes U)\subset\mathcal A,
          \qquad s\otimes T\longmapsto T(s),
\]
is an irreducible \(G_\infty\times G_f=G\)-representation.

The algebra spanned by the irreducible unitary action on \(S\) is
\(\operatorname{End}_{\mathbf C}(S)\), by the finite-dimensional
matrix-algebra argument written in PKT10.  Therefore an invariant
subspace of \(S\otimes U\) is of the form \(S\otimes U'\): in a
basis of \(S\), the matrix units extract each coefficient of a tensor
and insert it in every coordinate.  Stability under \(G_f\) makes
\(U'\) invariant, so irreducibility of \(U\) gives either zero or the
whole tensor product.  The evaluation map is equivariant and nonzero
because a nonzero map \(T\in U\) has some \(T(s)\ne0\).  Its kernel
is therefore zero, and its image is irreducible.  Each element of that
image is still an actual automorphic function.  This proves all fields
of the original \(\mathrm{IsAutRepU}\) for the image, and proves the
claimed algebraic decomposition.  In particular no Hilbert-space closure
is substituted for the required finite sum.
% END ADP4

% BEGIN ADP5
\paragraph{ADP5. Finite globally equivariant projections on a prescribed finite space.}
\label{claim:ADP5}
For every finite-dimensional subspace \(M\subset\mathcal A\), there
are finitely many original automorphic representations \(\sigma_i\)
and \(G\)-equivariant linear maps \(P_i:\mathcal A\to\mathcal A\)
such that \(P_i\mathcal A\subset\sigma_i\) and
\(\sum_iP_i f=f\) for all \(f\in M\).

ADP4 says that \(\mathcal A\), as a module over \(\mathbf C[G]\),
is a sum of simple submodules.  Choose a maximal direct family of simple
submodules by the usual union-of-chains argument.  If its sum were
proper, a simple submodule from the spanning family would not be
contained in it; its intersection with the sum is then zero, contradicting
maximality.  Hence this family is a direct-sum decomposition of
\(\mathcal A\).  Coordinate projections are \(G\)-equivariant.
Only finitely many coordinates occur in a finite basis of \(M\).
Select precisely those coordinates.  Their projections have the required
properties on every vector of \(M\).  Multiplicities cause no difficulty:
distinct summands may be isomorphic representations, and are kept as
distinct subspaces of the original function space.
% END ADP5

\subsection{The exact coefficient-to-function dictionary}

% BEGIN ADP6
\paragraph{ADP6. The finite projection and the prescribed infinity representation.}
\label{claim:ADP6}
The map \(p\) is a homomorphism, is surjective already on \(G_f\),
and has kernel \(G_\infty\).  The representation \(R\) is trivial on
\(G_f\), has continuous complex matrix coefficients, and satisfies
\[
 R(\gamma)E(v)=E\bigl(\xi(p(\gamma)_\ell)v\bigr)
                       \quad(\gamma\in\Gamma).
\]

The splitting of all selected places above \(\ell\), the exact local
sections, and the independent finite-coordinate adjustment are FNL4--FNL6.
They give surjectivity without changing the away coordinates.  If all
displayed coordinates of \(p(g)\) are one, every other finite coordinate
is one as well: at a selected split pair the opposite coordinate is the
inverse conjugate transpose of the displayed one.  Thus the kernel is
exactly \(G_\infty\).

Each tensor factor of \(R\) is the specified polynomial representation
evaluated on the complex infinite coordinate, so it is trivial on
finite-only elements.  Its matrix coefficients are polynomials in matrix
entries times an integral power of the nonzero determinant; they are
continuous, including negative determinant powers.  Finite tensor products
preserve this assertion.  The complex embedding of \(L\) indexing a
coordinate is fixed; this argument never regards \(\iota\) as a
continuous map on \(\overline{\mathbf Q}_\ell\).

For a rational matrix \(\gamma\), applying its prescribed embedding
to each entry of the selected local component is exactly its infinite
component at that embedding.  Apply the same polynomial representation
and tensor the equalities.  The original \(\mathrm{IsBasis}\) identity
for \(E,\xi\) gives the displayed rational compatibility.  This is the
literal dictionary for \(R,p,E\).
% END ADP6

% BEGIN ADP7
\paragraph{ADP7. Scalarization, its exact translation rule, and automorphy.}
\label{claim:ADP7}
For every coefficient function \(F\), the map \(\mathcal T_F\) of
\eqref{eq:adp-toform} is linear in \(F\) and in \(\lambda^*\), is
injective as a map of coefficient functions, and obeys
\[
 R_x\mathcal T_F(\lambda^*)
  =\mathcal T_{F(\,\cdot\,p(x))}
                    (\lambda^*\circ R(x^{-1})).
\]
If \(F\) has the original rational equivariance and an open finite
stabilizer, then every \(\mathcal T_F(\lambda^*)\) belongs to
\(\mathcal A\).

Linearity follows directly from the formula.  If every scalarization
vanishes, surjectivity of \(p\), invertibility of \(R(g)\), and
separation by linear functionals on finite-dimensional \(D\) give
\(E(F(y))=0\) for every \(y\), hence \(F=0\).  The translation
identity is the equation
\(R((gx)^{-1})=R(x^{-1})R(g^{-1})\), with the order as displayed.

By ADP6, rational equivariance says
\(E(F(p(\gamma)y))=R(\gamma)E(F(y))\).  Substitution proves left
\(\Gamma\)-invariance of the scalarization.  Right translation by
an open finite stabilizer fixes it, since \(R\) is trivial on \(G_f\).
Its \(G_\infty\)-translates lie in the image of the finite-dimensional
space \(D^*\), proving infinity finiteness.

For continuity, use the continuous finite projection
\(G=G_f\times G_\infty\to G_f\) from FNL2.  An open finite
stabilizer fixes \(F(p(g))\) locally in \(g\), while the dependence
on the infinite component is a continuous matrix coefficient of \(R\)
by ADP6.  Locally the scalarization is therefore a continuous matrix
coefficient with a fixed vector.  These are precisely all the original
\(\mathrm{IsAutFormU}\) conditions.
% END ADP7

% BEGIN ADP8
\paragraph{ADP8. Inverse scalarization on the original functions.}
\label{claim:ADP8}
Let \(\Phi:D^*\to\mathcal A\) be linear and suppose that
\[
 R_z\Phi(\lambda^*)
      =\Phi(\lambda^*\circ R(z^{-1}))\quad(z\in G_\infty).
\]
There exists a unique coefficient function \(F\) with the original
rational equivariance and \(\mathcal T_F=\Phi\).

Finite-dimensional biduality gives a unique \(\Psi(g)\in D\) with
\(\lambda^*(\Psi(g))=\Phi(\lambda^*)(g)\) for every functional.
The infinity relation is exactly
\(\Psi(gz)=R(z^{-1})\Psi(g)\).  Hence \(R(g)\Psi(g)\) is constant
on each fiber of \(p\), by ADP6.  Define
\[
                    F(p(g))=E^{-1}(R(g)\Psi(g)).
\]
This is independent of the lift \(g\).  Left rational invariance of
\(\Phi\) gives \(\Psi(\gamma g)=\Psi(g)\), hence
\(E(F(p(\gamma)y))=R(\gamma)E(F(y))\).  This is the exact rational
coefficient condition by ADP6.  Substitution gives
\(\mathcal T_F(\lambda^*)(g)=\Phi(\lambda^*)(g)\).  Uniqueness
follows from ADP7.  No coefficient field, basis equivalence, or
algebraic representation was changed in this inversion.
% END ADP8

% BEGIN ADP9
\paragraph{ADP9. Full original \texttt{formU}.}
\label{claim:ADP9}
The five clauses of the original \texttt{formU} all hold for the given
\(E,\xi,\iota\) and original projection \(p\): injectivity on the
rationally equivariant functions, automorphy under the original open
finite-stabilizer hypothesis, the exact right-translation identity,
surjectivity of the finite projection, and inverse scalarization for every
equivariant \(\Phi\).

ADP6 gives the finite projection and the exact rational dictionary.
ADP7 gives the first three function assertions, and ADP8 gives the final
surjectivity assertion, including the original function space in which
the inverse lies.  The argument works for any supplied finite-dimensional
basis equivalence, including a zero-dimensional coefficient space, and
uses no openness condition on \(\tau_0\) for this endpoint.  That
condition enters only when the level laws of a vector are used to produce
an open scalar-function stabilizer in ADP10.
% END ADP9

\subsection{Transport of finite automorphic projections}

% BEGIN ADP10
\paragraph{ADP10. The actual finite-level laws give simultaneous smoothness.}
\label{claim:ADP10}
Assume \(\ker\tau_0\) is open, and let \(F=(F_i)\in H.\mathrm{Big}\)
be fixed at one original away level \(K\) and one original Iwahori
depth \(b\).  Every scalarization of every \(F_i\) is an element of
\(\mathcal A\).

Choose a sufficiently small finite adelic neighborhood of one satisfying
the following finitely many conditions: its selected \(\ell\)-part lies
in \(I_1(b)\); its away part has the original level expression; and
its type coordinates lie in \(\ker\tau_0\).  Such a neighborhood
exists at every stated depth: the local principal congruence subgroup
\[
 \ker\bigl(\operatorname{GL}_n(O_{F_v})\longrightarrow
              \operatorname{GL}_n(O_{F_v}/\mathfrak m_{F_v}^{b+1})\bigr)
                         \ \subset I_1(b)
\]
is open.  To check the inclusion, eliminate a matrix congruent to the
identity modulo \(\mathfrak m_{F_v}^{b+1}\) into upper-unipotent,
diagonal and lower-unipotent factors.  Each pivot is congruent to one
and hence is a unit; each subtraction and division preserves the same
congruence.  The three factors therefore lie in the original
\(N\), \(T_c(b)\), and \(N^-(b)\), respectively.  The defining
Iwahori factorization puts their product in \(I_1(b)\).  The finite
product over the selected places is open as well.  Combine this with
FNL8--FNL9, continuity of the finite projections, openness of
the given \(K_q\)'s, and the assumed open kernel.  Outside these
finitely many places impose the usual integral conditions.  Thus every
element \(k\) of this neighborhood has
\(p(k)=(u,b')\), with \(u\in H.\mathrm{lev}(K)\),
\(\tau_C(u)=1\), and \(b'\in I_1(b)\).

The vector-valued level equation
\(\tau_C(u)F(y(u,1))=F(y)\) now says that every coordinate \(F_i\)
is fixed by this away translation.  Its \(I_1(b)\)-fixedness supplies
the local translation as well, so \(F_i(yp(k))=F_i(y)\).
Each scalar coordinate therefore has an open finite stabilizer.  Its
rational coefficient condition comes from \(H.\mathrm{Big}\), and
ADP7 applies.  This establishes a common neighborhood before the
coordinate or dual functional is selected; it does not assume separate
coordinate smoothness from vector-valued level invariance alone.
% END ADP10

% BEGIN ADP11
\paragraph{ADP11. Equivariant projections preserve the exact coefficient equations.}
\label{claim:ADP11}
For \(F\) as in ADP10, there are finitely many original automorphic
representations \(\sigma_t\) and vectors \(G_t\in H.\mathrm{Big}\)
such that
\[
 F=\sum_tG_t,\qquad
            \mathcal T_{(G_t)_i}(\lambda^*)\in\sigma_t
                         \quad\text{for every }i,\lambda^*.
\]

Apply ADP5 to the finite-dimensional span of all
\(\mathcal T_{F_i}(\lambda^*)\), as \(i\) ranges over its finite
index set and \(\lambda^*\) over a basis of \(D^*\).  Obtain the
global equivariant projections \(P_t\).  Since \(P_t\) commutes with
the infinity action, \(P_t\mathcal T_{F_i}\) satisfies the exact
equivariance required in ADP8.  Let \((G_t)_i\) be its inverse
scalarization.  These functions have rational coefficient equivariance
and have all scalarizations in \(\sigma_t\).

It remains to check the base away coefficient equation defining
\(H.\mathrm{Big}\).  For \(u\) in its original base subgroup, choose
\(k_u\in G_f\) with \(p(k_u)=(u,1)\), using ADP6.  ADP7 gives
\[
 \mathcal T_{F_i(\,\cdot\,(u,1))}(\lambda^*)
      =R_{k_u}\mathcal T_{F_i}(\lambda^*).
\]
The original matrix equation is a finite linear combination of these
identities with coefficients \((\tau_C(u))_{ij}\).  Apply \(P_t\),
commute it past \(R_{k_u}\), and use its complex linearity.  ADP7
injectivity recovers the identical matrix equation for \(G_t\).
Thus \(G_t\in H.\mathrm{Big}\).  Finally
\(\sum_tP_t\mathcal T_{F_i}=\mathcal T_{F_i}\); inverse
scalarization gives \(\sum_tG_t=F\).  This checks the actual
coefficient action, including all matrix mixing, rather than replacing
it by a scalar invariance assumption.
% END ADP11

% BEGIN ADP12
\paragraph{ADP12. Full original \texttt{decompU} and its exact remaining boundary.}
\label{claim:ADP12}
For every datum and vector in the original \texttt{decompU} statement,
\[
 F\in\sum_{\sigma:\,\mathrm{IsAutRepU}(\sigma)}
       Y_{\sigma},\qquad
 Y_{\sigma}=H.\mathrm{Big}\cap
       \operatorname{span}_{\mathbf C}
       \{G:\mathcal T_{G_i}(\lambda^*)\in\sigma
                                   \text{ for all }i,\lambda^*\}.
\]
This is the literal source \(YOf\), with the specified \(E,\xi\).

ADP10 proves automorphy for all scalarizations from the exact hypotheses
of \texttt{decompU}.  ADP11 gives a finite sum of vectors in the
displayed spaces.  Each vector lies even in the generating set before
taking its span.  This proves the endpoint with the same functions,
coefficient representation, embeddings, and rational group.

The theorem makes no arithmetic assertion about \(\sigma\).  To turn
\(Y_\sigma\ne0\) into a \(\mathrm{PartAltAlt}\) one must still
construct its Galois representation, spherical polynomial readings,
local ordinary comparison, auxiliary local correspondence, and the
original point-property record.  The source's
\texttt{partU\_alt}/\texttt{partX\_alt} is the separate producer of
those data.  Once it is supplied, discard the zero \(Y_\sigma\)'s,
decorate each remaining one without changing its subspace, and apply
CFD12.  No assertion of arithmetic nonvanishing or completion of the
same-ideal Hida theorem follows solely from this decomposition.
% END ADP12
